% ===== PRÉAMBULE CANONIQUE — à coller VERBATIM en tête de CHAQUE .tex =====
\documentclass[11pt,a4paper]{article}
\usepackage[utf8]{inputenc}\usepackage[T1]{fontenc}\usepackage{lmodern}
\usepackage[french]{babel}
\usepackage{amsmath,amssymb,mathtools}\usepackage{icomma}
\usepackage[version=4]{mhchem}          % \ce{...} : équations & formules
\usepackage{siunitx}\sisetup{locale=FR} % \qty{}{}, \si{}, \ang{}
\usepackage{chemfig}                    % \chemfig{...} : structures développées
\usepackage{xcolor}\usepackage{colortbl}
\usepackage{tikz}\usepackage{pgfplots}\pgfplotsset{compat=1.18}
\usepackage[most]{tcolorbox}\usepackage{enumitem}\usepackage{array,booktabs}
\usepackage[a4paper,margin=2.2cm]{geometry}
\usetikzlibrary{arrows.meta,calc,angles,quotes,positioning,patterns,backgrounds,shapes.geometric}
\definecolor{primary}{RGB}{222,49,99}\definecolor{primarydark}{RGB}{150,22,55}
\definecolor{defcol}{RGB}{0,114,178}\definecolor{methcol}{RGB}{52,92,148}
\definecolor{excol}{RGB}{0,135,90}\definecolor{attncol}{RGB}{200,40,55}
\tcbset{boxbase/.style={enhanced,breakable,boxrule=0.9pt,arc=2.5pt,
  left=3mm,right=3mm,top=2.3mm,bottom=2.3mm,fonttitle=\bfseries,coltitle=white}}
% --- boîtes TUTORIEL (noms EXACTS, ne pas renommer) ---
\newtcolorbox{cle}[1][]{boxbase,colback=defcol!6,colframe=defcol,
  title={\textsc{À retenir}\ifx\relax#1\relax\relax\else\textmd{ : \itshape #1}\fi}}
\newtcolorbox{methode}[1][]{boxbase,colback=methcol!7,colframe=methcol,sharp corners,
  title={\textsc{Méthode}\ifx\relax#1\relax\relax\else\textmd{ --- \itshape #1}\fi}}
\newtcolorbox{exemplebox}[1][]{boxbase,colback=excol!5,colframe=excol,
  title={\textsc{Exemple}\ifx\relax#1\relax\relax\else\textmd{ : \itshape #1}\fi}}
\newtcolorbox{piege}[1][]{boxbase,colback=attncol!6,colframe=attncol,
  title={\textsc{Piège}\ifx\relax#1\relax\relax\else\textmd{ : \itshape #1}\fi}}
\newtcolorbox{remarque}[1][]{boxbase,colback=black!4,colframe=black!45,
  title={\textsc{Remarque}\ifx\relax#1\relax\relax\else\textmd{ : \itshape #1}\fi}}
% --- boîtes EXERCICE / CORRIGÉ (noms EXACTS, auto-comptées) ---
\newtcolorbox[auto counter]{exo}[1][]{boxbase,colback=primary!4,colframe=primary,
  title={\textsc{Exercice}~\thetcbcounter\ifx\relax#1\relax\relax\else\textmd{ --- \itshape #1}\fi}}
\newtcolorbox[auto counter]{corrige}[1][]{boxbase,colback=excol!4,colframe=excol,
  title={\textsc{Corrigé}~\thetcbcounter\ifx\relax#1\relax\relax\else\textmd{ --- \itshape #1}\fi}}
\newcommand{\R}{\mathbb{R}}\newcommand{\N}{\mathbb{N}}\newcommand{\Z}{\mathbb{Z}}\newcommand{\ds}{\displaystyle}
% (un \usepackage{fancyhdr}+en-tête est possible ; il est ignoré côté web)
% =========================================================================
\begin{document}
\begin{center}\color{primarydark}\large\bfseries Thème 4 — Corrigé \quad$\bullet$\quad Niveau Facile\end{center}

\begin{corrige}[les réflexes de calcul]
\begin{enumerate}[label=\alph*)]
  \item $\log(100)=2$ ; $\log(10)=1$ ; $\log(1)=0$ ; $\log(0{,}1)=-1$ ; $\log(0{,}01)=-2$.
  \item $\dfrac{0{,}06}{1}=0{,}06$ ; $\dfrac{0{,}06}{2}=0{,}03$ ; $\dfrac{0{,}06}{3}=0{,}02$.
\end{enumerate}
\end{corrige}

\begin{corrige}[quand le potentiel vaut $E_{\mathrm{r}}^{\circ}$]
\begin{enumerate}[label=\alph*)]
  \item $\dfrac{[\ce{Fe^{3+}}]}{[\ce{Fe^{2+}}]}=1$, donc $\log(1)=0$ :
        $E=0{,}77+0{,}06\times 0=+0{,}77$~V.
  \item $[\ce{Cu^{2+}}]=1$, le métal a une activité $1$ : $\log(1)=0$, donc
        $E=0{,}34+0=+0{,}34$~V. En conditions standard, $E=E_{\mathrm{r}}^{\circ}$.
\end{enumerate}
\end{corrige}

\begin{corrige}[électrode diluée]
$[\mathrm{Red}]=1$ (métal solide), $[\mathrm{Ox}]=[\ce{Cu^{2+}}]=0{,}01$, $n=2$ :
\[
E=0{,}34+\frac{0{,}06}{2}\log(0{,}01)=0{,}34+0{,}03\times(-2)=0{,}34-0{,}06=+0{,}28~\text{V}.
\]
\end{corrige}

\begin{corrige}[force électromotrice standard]
\begin{enumerate}[label=\alph*)]
  \item Le couple de $E_{\mathrm{r}}^{\circ}$ le plus élevé est la \textbf{cathode} :
        \ce{Cu^{2+}}/\ce{Cu} ($+0{,}34$). L'\textbf{anode} est \ce{Zn^{2+}}/\ce{Zn} ($-0{,}76$).
  \item $\Delta E^{\circ}=0{,}34-(-0{,}76)=+1{,}10$~V.
\end{enumerate}
\end{corrige}

\begin{corrige}[cathode, anode, f.é.m.]
\begin{enumerate}[label=\alph*)]
  \item Cathode : \ce{Ag+}/\ce{Ag} ($+0{,}80$, le plus haut) ; anode : \ce{Fe^{2+}}/\ce{Fe}
        ($-0{,}44$).
  \item $\Delta E^{\circ}=0{,}80-(-0{,}44)=+1{,}24$~V.
\end{enumerate}
\end{corrige}

\end{document}
